\documentclass[14pt,a4paper]{extarticle}
\usepackage{iftex}
\ifPDFTeX
    \usepackage[T2A]{fontenc}
    \usepackage[utf8]{inputenc}
    \usepackage[russian]{babel}
\else
    \usepackage{fontspec}
    \usepackage[russian]{babel}
\fi
\usepackage{geometry}
\usepackage{setspace}
\usepackage{indentfirst}
\usepackage{titlesec}
\usepackage{longtable}
\usepackage{array}
\usepackage{hyperref}
\usepackage{amsmath}
\usepackage{xcolor}
\usepackage{listings}
\usepackage{graphicx}
\usepackage{tikz}
\usetikzlibrary{positioning,arrows.meta}

\geometry{left=25mm,right=15mm,top=20mm,bottom=20mm}
\onehalfspacing
\setlength{\parindent}{1.25cm}
\setlength{\parskip}{0pt}

\titleformat{\section}{\normalfont\bfseries}{\thesection.}{0.5em}{}
\titleformat{\subsection}{\normalfont\bfseries}{\thesubsection.}{0.5em}{}
\titleformat{\subsubsection}{\normalfont\bfseries}{\thesubsubsection.}{0.5em}{}

\hypersetup{
    colorlinks=true,
    linkcolor=black,
    urlcolor=black,
    citecolor=black
}

\lstdefinestyle{code}{
    basicstyle=\small\ttfamily,
    breaklines=true,
    tabsize=4,
    showstringspaces=false,
    frame=single,
    rulecolor=\color{black!20},
    keywordstyle=\color{blue!60!black},
    commentstyle=\color{black!55},
    stringstyle=\color{teal!60!black}
}

\lstset{style=code}

\begin{document}

\begin{center}
МИНИСТЕРСТВО НАУКИ И ВЫСШЕГО ОБРАЗОВАНИЯ РФ\\
ФЕДЕРАЛЬНОЕ ГОСУДАРСТВЕННОЕ БЮДЖЕТНОЕ\\
ОБРАЗОВАТЕЛЬНОЕ УЧРЕЖДЕНИЕ ВЫСШЕГО ОБРАЗОВАНИЯ\\
«ЛИПЕЦКИЙ ГОСУДАРСТВЕННЫЙ ТЕХНИЧЕСКИЙ УНИВЕРСИТЕТ»\\
\vspace{0.5cm}
Институт компьютерных наук\\
Кафедра прикладной математики
\vspace{2cm}

\textbf{КУРСОВАЯ РАБОТА}\\
по дисциплине: «Современные платформы разработки программного обеспечения»\\
на тему: «Разработка бухгалтерского приложения на C\# и WPF»
\vspace{2.5cm}

\begin{flushright}
Студент \underline{\hspace{4cm}} \\
группа \underline{\hspace{3cm}} \\
\vspace{0.5cm}
Руководитель \underline{\hspace{4cm}}
\end{flushright}

\vfill
Липецк \the\year{} г.
\end{center}

\newpage
\tableofcontents
\newpage

\section{Техническое задание}

\subsection{Введение}
Наименование программы: \textit{BuhUchet}.\\
Тип программы: настольное бухгалтерское приложение.\\
Платформа: Windows, .NET 8, WPF.

Программа предназначена для ведения журнала хозяйственных операций, работы со справочником контрагентов, просмотра плана счетов и формирования оборотно-сальдовой ведомости (ОСВ).

\subsection{Основание для разработки}
Разработка выполнена в рамках учебного проекта по дисциплине «Современные платформы разработки программного обеспечения».

\subsection{Назначение разработки}
Основное назначение:
\begin{itemize}
    \item регистрация и редактирование бухгалтерских проводок;
    \item хранение справочников (контрагенты, план счетов);
    \item построение отчета ОСВ за заданный период;
    \item экспорт и импорт отчетных данных в формате JSON.
\end{itemize}

\subsection{Требования к программе}
\textbf{Средства разработки:}
\begin{itemize}
    \item язык программирования: C\#;
    \item UI-фреймворк: WPF;
    \item целевая платформа: \texttt{net8.0-windows};
    \item библиотека сериализации: Newtonsoft.Json.
\end{itemize}

\textbf{Функциональные модули:}
\begin{itemize}
    \item модуль журнала операций;
    \item модуль контрагентов;
    \item модуль плана счетов;
    \item модуль отчетности (ОСВ);
    \item модуль автосохранения и хранения данных;
    \item модуль автодополнения полей ввода.
\end{itemize}

\subsection{Функциональные требования}
\begin{itemize}
    \item добавление, изменение и удаление записей в журнале;
    \item фильтрация журнала по периоду и строке поиска;
    \item хранение ИНН и банковских реквизитов контрагентов;
    \item расчет входящих, оборотных и исходящих показателей ОСВ;
    \item экспорт ОСВ в JSON и загрузка ранее сохраненного отчета;
    \item автоматическое сохранение рабочих данных в \texttt{data.json}.
\end{itemize}

\subsection{Нефункциональные требования}
\begin{itemize}
    \item удобный графический интерфейс;
    \item устойчивость к отсутствию или повреждению файла данных;
    \item минимальное время отклика при работе с таблицами;
    \item корректная работа с русской локалью при вводе сумм.
\end{itemize}

\subsection{Требования к данным}
Основной файл данных: \texttt{data.json}.\\
Структура данных:
\begin{itemize}
    \item \texttt{Journal} --- список проводок;
    \item \texttt{Counterparties} --- список контрагентов;
    \item \texttt{Accounts} --- план счетов;
    \item \texttt{SavedAt} --- дата и время последнего сохранения.
\end{itemize}

\subsection{Этапы разработки}
\begin{itemize}
    \item проектирование модели данных;
    \item реализация интерфейса WPF;
    \item реализация бизнес-логики (журнал, контрагенты, ОСВ);
    \item добавление автодополнения и автосохранения;
    \item тестирование основных пользовательских сценариев.
\end{itemize}

\section{Описание структуры приложения}

\subsection{Диаграмма классов}
Ниже приведена UML-подобная диаграмма ключевых классов приложения.

\begin{center}
\begin{tikzpicture}[
    box/.style={draw, rounded corners=3pt, align=left, font=\footnotesize, text width=6.0cm, inner sep=7pt},
    smallbox/.style={draw, rounded corners=3pt, align=left, font=\footnotesize, text width=5.8cm, inner sep=7pt},
    rel/.style={-{Latex[length=2.6mm]}, very thick}
]
\node[box] (mw) at (0,0) {\textbf{MainWindow}\\
Поля: \_journal, \_counterparties, \_accounts\\
Методы: LoadData(), TrySaveData(), BtnBuildOsv\_Click()};

\node[smallbox, below=1.3cm of mw] (ds) {\textbf{DataService}\\
BuildOsv(journal, accounts, from, to): List\textless OsvRow\textgreater\\
FilterByCounterparty(...): List\textless JournalEntry\textgreater};

\node[smallbox, below=1.0cm of ds] (ac) {\textbf{AutoCompleteService}\\
Attach(textBox, sourceProvider, onSelected)\\
GetMatches(), ShowSuggestions(), HandleTab()};

\node[smallbox, below=1.0cm of ac] (st) {\textbf{AppDataStore}\\
LoadOrCreate(path): SaveFile\\
Save(path, data): void; DefaultPath: string};

\node[box, right=2.8cm of mw] (sf) {\textbf{SaveFile}\\
Journal: List\textless JournalEntry\textgreater\\
Counterparties: List\textless Counterparty\textgreater\\
Accounts: List\textless AccountPlan\textgreater\\
SavedAt: DateTime};

\node[smallbox, below=1.0cm of sf] (je) {\textbf{JournalEntry}\\
Id, Date, DebitAccount, CreditAccount, Amount\\
Counterparty, CounterpartyInn, CounterpartyBankAccount};

\node[smallbox, below=1.0cm of je] (cp) {\textbf{Counterparty}\\
Name, Inn, BankAccount, PersonalAccount, Type};

\node[smallbox, below=1.0cm of cp] (ap) {\textbf{AccountPlan}\\
Code, Name, Type};

\node[smallbox, below=1.0cm of ap] (orow) {\textbf{OsvRow}\\
Account, AccountName\\
OpenDebit, OpenCredit, TurnDebit, TurnCredit\\
CloseDebit, CloseCredit};

\draw[rel] (mw.east) -- node[above, font=\footnotesize] {использует} (sf.west);
\draw[rel] (mw.south east) .. controls +(1.2,-0.9) and +(-1.2,0.8) .. node[right, font=\footnotesize] {вызывает} (ds.west);
\draw[rel] (mw.south) -- node[left, font=\footnotesize] {использует} (ac.north);
\draw[rel] (mw.south) .. controls +(0.1,-4.2) and +(-1.1,0.9) .. node[left, font=\footnotesize] {чтение/запись} (st.north);

\draw[rel] (sf.south) -- node[right, font=\footnotesize] {содержит} (je.north);
\draw[rel] (sf.south west) .. controls +(-0.3,-1.5) and +(0.9,1.2) .. node[right, font=\footnotesize] {содержит} (cp.north);
\draw[rel] (sf.south west) .. controls +(-0.8,-2.8) and +(0.9,1.2) .. node[right, font=\footnotesize] {содержит} (ap.north);

\draw[rel] (ds.east) .. controls +(1.2,0.0) and +(-1.2,0.3) .. node[below, font=\footnotesize] {формирует} (orow.west);
\draw[rel] (ds.east) .. controls +(1.8,0.9) and +(-1.2,-0.5) .. node[above, font=\footnotesize] {читает} (je.west);
\draw[rel] (ds.east) .. controls +(2.0,0.2) and +(-1.2,0.0) .. node[above, font=\footnotesize] {читает} (ap.west);
\end{tikzpicture}
\end{center}

\subsection{Описание классов}
\subsubsection{Назначение классов предметной области}
Предметная область приложения представлена сущностями:
\begin{itemize}
    \item \texttt{JournalEntry} --- бухгалтерская проводка (журнальная запись);
    \item \texttt{Counterparty} --- контрагент (справочник контрагентов);
    \item \texttt{AccountPlan} --- элемент плана счетов;
    \item \texttt{OsvRow} --- строка оборотно-сальдовой ведомости;
    \item \texttt{SaveFile} --- корневая структура сохранения данных.
\end{itemize}

\subsubsection{Назначение сервисных классов}
Сервисные компоненты:
\begin{itemize}
    \item \texttt{AppDataStore} --- чтение и запись файла данных \texttt{data.json};
    \item \texttt{DataService} --- расчёт ОСВ за период и фильтрация операций;
    \item \texttt{AutoCompleteService} --- автодополнение (Tab и всплывающий список) в полях ввода.
\end{itemize}

\subsection{Таблица описания классов}
Основные файлы:
\begin{itemize}
    \item \texttt{MainWindow.xaml} --- разметка главного окна;
    \item \texttt{MainWindow.xaml.cs} --- обработчики UI и сценарии работы;
    \item \texttt{Models.cs} --- классы предметной области;
    \item \texttt{DataService.cs} --- сервис расчетов отчетов;
    \item \texttt{AutoCompleteService.cs} --- автодополнение в полях ввода;
    \item \texttt{AppDataStore.cs} --- загрузка и сохранение данных.
\end{itemize}

\subsection{Описание ключевых классов}
\begin{longtable}{|p{3.6cm}|p{4.4cm}|p{6.8cm}|}
\hline
\textbf{Класс} & \textbf{Назначение} & \textbf{Ключевые элементы} \\
\hline
\texttt{MainWindow} & Главный контроллер интерфейса & Навигация по разделам, обработчики кнопок, редактирование записей, экспорт/импорт ОСВ \\
\hline
\texttt{JournalEntry} & Модель проводки & Id, Date, DebitAccount, CreditAccount, Counterparty, Amount, Description \\
\hline
\texttt{Counterparty} & Модель контрагента & Name, Inn, BankAccount, PersonalAccount, Type \\
\hline
\texttt{AccountPlan} & Модель счета учета & Code, Name, Type \\
\hline
\texttt{OsvRow} & Строка оборотно-сальдовой ведомости & OpenDebit/OpenCredit, TurnDebit/TurnCredit, CloseDebit/CloseCredit \\
\hline
\texttt{DataService} & Формирование отчетов & \texttt{BuildOsv()}, \texttt{FilterByCounterparty()} \\
\hline
\texttt{AutoCompleteService} & Подсказки при вводе & \texttt{Attach()}, обработка Tab/Escape/Down, popup-список \\
\hline
\texttt{AppDataStore} & Работа с файлом данных & \texttt{LoadOrCreate()}, \texttt{Save()}, путь \texttt{DefaultPath} \\
\hline
\end{longtable}

\subsection{Взаимодействие компонентов}
\begin{enumerate}
    \item При старте \texttt{MainWindow} загружает данные через \texttt{AppDataStore}.
    \item Данные привязываются к \texttt{DataGrid} в виде \texttt{ObservableCollection}.
    \item Пользователь изменяет записи; форма валидирует поля.
    \item После изменений запускается отложенное автосохранение.
    \item По запросу строится ОСВ через \texttt{DataService}.
    \item Отчет может быть выгружен в JSON и загружен обратно.
\end{enumerate}

\subsection{Архитектурные особенности}
\begin{itemize}
    \item разделение на слой UI, модели и сервисы;
    \item применение \texttt{INotifyPropertyChanged} для реактивного обновления интерфейса;
    \item использование \texttt{DispatcherTimer} для debounce-автосохранения;
    \item работа с постоянным хранилищем через JSON-сериализацию.
\end{itemize}

\section{Руководство оператора}

\subsection{Общие сведения о программе}
Приложение содержит 4 раздела:
\begin{itemize}
    \item Журнал операций;
    \item Контрагенты;
    \item План счетов;
    \item Отчеты (ОСВ).
\end{itemize}

\subsection{Описание элементов интерфейса}
\textbf{Левая панель навигации} содержит кнопки разделов и отображение статуса файла данных (\texttt{data.json}).\\
\textbf{Журнал операций} включает:
\begin{itemize}
    \item таблицу операций (DataGrid) с колонками: №, Дата, Дт, Кт, Контрагент, Сумма, Содержание;
    \item панель фильтрации по периоду (DatePicker «с/по») и строку поиска;
    \item форму редактирования/добавления записи с валидацией обязательных полей;
    \item быстрые подсказки в полях (автодополнение по Tab и выпадающий список).
\end{itemize}
\textbf{Контрагенты} включают таблицу справочника и форму редактирования реквизитов (ИНН, банковский счёт, лицевой счёт, тип).\\
\textbf{План счетов} отображается в виде таблицы (код, наименование, тип).\\
\textbf{Отчёты (ОСВ)} включают выбор периода, таблицу ОСВ, строку итогов и кнопки экспорта/импорта JSON.

\subsection{Необходимое программное обеспечение}
\begin{itemize}
    \item ОС: Windows 10/11;
    \item .NET: целевая платформа \texttt{net8.0-windows};
    \item IDE для разработки/запуска: Visual Studio 2022+ (Desktop development with C\#);
    \item (опционально) LaTeX-дистрибутив для сборки отчёта: TeX Live / MiKTeX.
\end{itemize}

\subsection{Необходимое аппаратное обеспечение}
\begin{longtable}{|p{5.2cm}|p{4.8cm}|p{4.8cm}|}
\hline
\textbf{Параметр} & \textbf{Минимальные требования} & \textbf{Рекомендуемые требования} \\
\hline
Процессор & 2 ядра, 2.0 ГГц & 4 ядра, 2.5 ГГц+ \\
\hline
ОЗУ & 4 ГБ & 8 ГБ+ \\
\hline
Накопитель & 300 МБ свободно & 1 ГБ свободно \\
\hline
Экран & 1366x768 & 1920x1080 \\
\hline
\end{longtable}

\subsection{Выполнение программы}
\begin{enumerate}
    \item Открыть проект \texttt{BuhUchet.sln} в Visual Studio.
    \item Выполнить сборку и запуск (\texttt{F5}).
    \item После запуска перейти в нужный раздел через левое меню.
\end{enumerate}

\subsection{Входные и выходные данные}
\textbf{Входные данные:}
\begin{itemize}
    \item данные проводок (дата, счета, контрагент, сумма, описание);
    \item реквизиты контрагентов (ИНН, счета, тип);
    \item период отчёта ОСВ (даты «с/по»);
    \item JSON-файл отчёта ОСВ при импорте.
\end{itemize}
\textbf{Выходные данные:}
\begin{itemize}
    \item содержимое таблиц в UI;
    \item сформированная ОСВ (таблица и итоги);
    \item файл \texttt{data.json} (автосохранение);
    \item экспортированный JSON-файл отчёта ОСВ.
\end{itemize}

\subsection{Сообщения оператору}
Приложение использует строку статуса внизу окна и короткие сообщения:
\begin{itemize}
    \item информационные: «[OK] Операция сохранена», «[OK] Контрагент сохранён», «[OK] ОСВ сформирована: N счетов», «[OK] Экспортировано: имя\_файла»;
    \item предупреждения/ошибки: «[WARN] Укажите период», «[WARN] Введите корректную сумму», «[WARN] Ошибка экспорта: ...», «[WARN] Ошибка загрузки отчёта: ...», «[WARN] Ошибка автосохранения: ...».
\end{itemize}

\subsection{Примеры использования}
\textbf{Пример 1: добавление проводки.}
\begin{enumerate}
    \item Раздел «Журнал операций» \(\rightarrow\) кнопка «Добавить».
    \item Заполнить: Дата, Дт, Кт, Контрагент, Сумма, Содержание.
    \item Нажать «Сохранить». Внизу появится сообщение «[OK] Операция сохранена».
\end{enumerate}
\textbf{Пример 2: формирование ОСВ и экспорт.}
\begin{enumerate}
    \item Раздел «Отчёты (ОСВ)» \(\rightarrow\) выбрать период.
    \item Нажать «Сформировать». Проверить итоги.
    \item Нажать «Экспорт ОСВ (JSON)» и сохранить файл.
\end{enumerate}

\subsection{Работа с журналом операций}
\begin{enumerate}
    \item Нажать «Добавить».
    \item Заполнить дату, счета Дт/Кт, контрагента, сумму и описание.
    \item Нажать «Сохранить».
    \item Для изменения выбрать запись в таблице, отредактировать и сохранить.
    \item Для удаления выбранной записи нажать кнопку удаления и подтвердить действие.
\end{enumerate}

\subsection{Работа с контрагентами}
\begin{enumerate}
    \item Перейти в раздел «Контрагенты».
    \item Добавить или выбрать существующего контрагента.
    \item Заполнить реквизиты (ИНН, банковский счет, лицевой счет, тип).
    \item Сохранить изменения.
\end{enumerate}

\subsection{Формирование ОСВ}
\begin{enumerate}
    \item Перейти в раздел «Отчеты (ОСВ)».
    \item Указать период «с» и «по».
    \item Нажать «Сформировать».
    \item При необходимости выполнить экспорт отчета в JSON.
\end{enumerate}

\section{Описание контрольного примера}

\subsection{Назначение}
Контрольный пример демонстрирует корректность:
\begin{itemize}
    \item ввода и хранения проводок;
    \item расчета итогов по счетам;
    \item формирования ОСВ за выбранный период.
\end{itemize}

\subsection{Исходные данные}
Период: 01.01.2026--31.01.2026.\\
Проводки:
\begin{enumerate}
    \item Дт 10 Кт 60, сумма 10000;
    \item Дт 20 Кт 10, сумма 4000;
    \item Дт 62 Кт 90, сумма 15000;
    \item Дт 51 Кт 62, сумма 7000.
\end{enumerate}

\subsection{Дополнение для примера}
В файле \texttt{data.json} данные хранятся в виде объекта \texttt{SaveFile}. Для проверки также допустимо подготовить тестовый \texttt{data.json} вручную (например, добавить несколько счетов в \texttt{Accounts} и несколько операций в \texttt{Journal}).

\subsection{Сценарий для проверки}
\begin{enumerate}
    \item Внести проводки в журнал.
    \item Сформировать ОСВ за январь 2026.
    \item Проверить обороты по счетам 10, 20, 51, 60, 62, 90.
    \item Проверить корректность итоговых строк.
\end{enumerate}

\subsection{Визуализация}
Результатом выполнения сценария является заполненная таблица ОСВ в разделе «Отчёты (ОСВ)» и строка итогов. Таблица включает колонки начального сальдо, оборотов и конечного сальдо.

\subsection{Результаты расчёта}
\textbf{Ожидаемые обороты за период:}
\begin{itemize}
    \item счёт 10: оборот Дт = 10000, оборот Кт = 4000;
    \item счёт 20: оборот Дт = 4000;
    \item счёт 60: оборот Кт = 10000;
    \item счёт 62: оборот Дт = 15000, оборот Кт = 7000;
    \item счёт 90: оборот Кт = 15000;
    \item счёт 51: оборот Дт = 7000.
\end{itemize}
Конечные сальдо рассчитываются формулами, используемыми в \texttt{OsvRow}:
\begin{equation}
\text{CloseDebit}=\max(0,\;OpenDebit-OpenCredit+TurnDebit-TurnCredit)
\end{equation}
\begin{equation}
\text{CloseCredit}=\max(0,\;OpenCredit-OpenDebit+TurnCredit-TurnDebit)
\end{equation}

\subsection{Проверка программы}
\subsubsection{Состав технических средств}
ПК под управлением Windows 10/11, установленный .NET 8 runtime (или Visual Studio для запуска из IDE).

\subsubsection{Формирование исходных данных}
Исходные данные формируются через пользовательский интерфейс (журнал операций), либо путём заполнения \texttt{data.json} и последующего запуска приложения.

\subsubsection{Описание действий оператора}
\begin{enumerate}
    \item Запустить приложение.
    \item Внести операции (см. раздел 4.2).
    \item Перейти в «Отчёты (ОСВ)», указать период, сформировать ОСВ.
    \item Сравнить обороты и итоги со значениями из раздела 4.3.
\end{enumerate}

\subsection{Пример выполнения}
После ввода операций и формирования ОСВ оператор видит:
\begin{itemize}
    \item список задействованных счетов;
    \item обороты по дебету/кредиту;
    \item рассчитанные конечные сальдо;
    \item строку итогов (суммы оборотов и конечных сальдо).
\end{itemize}

\subsection{Проверка ошибок}
\begin{longtable}{|p{6.3cm}|p{4.7cm}|p{4.0cm}|}
\hline
\textbf{Сценарий} & \textbf{Ожидаемая реакция} & \textbf{Результат} \\
\hline
Пустая дата операции & Сообщение «Укажите дату операции» & Блокируется сохранение \\
\hline
Некорректная сумма (буквы/0/отрицательная) & Сообщение «Введите корректную сумму» & Блокируется сохранение \\
\hline
Экспорт в недоступную папку & «Ошибка экспорта: ...» & Сообщение об ошибке \\
\hline
Импорт повреждённого JSON & «Ошибка загрузки отчёта: ...» & Сообщение об ошибке \\
\hline
\end{longtable}
\begin{itemize}
    \item для каждого использованного счета отображаются обороты по дебету и кредиту;
    \item рассчитаны начальные и конечные сальдо;
    \item итоговые суммы по ОСВ отображаются в статусной области отчета;
    \item экспортированный JSON содержит разделы \texttt{Rows} и \texttt{Totals}.
\end{itemize}

\section{Сообщения об ошибках}

\subsection{Классификация ошибок}
\begin{itemize}
    \item ошибки валидации пользовательского ввода;
    \item ошибки чтения/записи файла;
    \item ошибки парсинга формата отчета.
\end{itemize}

\subsection{Таблица ошибок}
\begin{longtable}{|p{2.4cm}|p{5.1cm}|p{4.0cm}|p{3.1cm}|}
\hline
\textbf{Код} & \textbf{Ситуация} & \textbf{Сообщение} & \textbf{Действие} \\
\hline
VAL-001 & Не указана дата операции & «Укажите дату операции» & Блокировка сохранения \\
\hline
VAL-002 & Некорректная сумма & «Введите корректную сумму» & Блокировка сохранения \\
\hline
IO-001 & Ошибка экспорта отчета & «Ошибка экспорта: ...» & Показ статуса ошибки \\
\hline
IO-002 & Ошибка загрузки отчета & «Ошибка загрузки отчёта: ...» & Показ статуса ошибки \\
\hline
IO-003 & Ошибка автосохранения & «Ошибка автосохранения: ...» & Уведомление оператора \\
\hline
\end{longtable}

\subsection{Примеры обработки ошибок в коде}
\textbf{Валидация полей при сохранении проводки} (пример логики из \texttt{MainWindow.xaml.cs}):
\begin{lstlisting}[language={[Sharp]C}]
if (PickDate.SelectedDate == null) { ShowStatus("Укажите дату операции", true); return; }
if (string.IsNullOrWhiteSpace(TxtDebit.Text)) { ShowStatus("Укажите счёт дебета", true); return; }
if (string.IsNullOrWhiteSpace(TxtCredit.Text)) { ShowStatus("Укажите счёт кредита", true); return; }
if (!TryParseAmount(TxtAmount.Text, out decimal amount) || amount <= 0)
{ ShowStatus("Введите корректную сумму", true); return; }
\end{lstlisting}

\textbf{Обработка исключений при экспорте/импорте отчёта} реализована через \texttt{try/catch} с выводом сообщения в строку статуса.

\subsection{Визуальное отображение ошибок}
Ошибки отображаются в нижней статусной строке:
\begin{itemize}
    \item ошибки: префикс «[WARN]», цвет текста --- \textit{LightSalmon};
    \item успешные операции: префикс «[OK]», цвет текста --- \textit{LightGreen}.
\end{itemize}

\subsection{Логирование ошибок}
В текущей версии приложения отдельный файл логов не ведётся: информация выводится в UI (строка статуса) и в исключительных ситуациях может быть доступна через стандартные средства отладки.\\
\textbf{Рекомендуемое улучшение}: добавление записи ошибок в файл (например, \texttt{error\_log.txt}) в формате \texttt{[Дата] [Код] [Сообщение]} в обработчиках исключений автосохранения, экспорта и импорта.

\subsection{Рекомендации для тестировщиков}
\begin{itemize}
    \item проверить валидацию обязательных полей (пустая дата, пустые счета, неверная сумма);
    \item проверить корректность расчёта ОСВ на небольшом наборе проводок с ручным подсчётом;
    \item проверить автосохранение: изменить данные и убедиться, что \texttt{data.json} обновляется;
    \item проверить импорт ОСВ из корректного JSON и из повреждённого JSON.
\end{itemize}

\section{Список литературы}
\begin{enumerate}
    \item ГОСТ 19.401-78. ЕСПД. Текст программы. Требования к содержанию и оформлению.
    \item ГОСТ 19.505-79. ЕСПД. Руководство оператора. Требования к содержанию и оформлению.
    \item Microsoft Docs. WPF documentation. \url{https://learn.microsoft.com/dotnet/desktop/wpf/}
    \item Microsoft Docs. C\# documentation. \url{https://learn.microsoft.com/dotnet/csharp/}
    \item Newtonsoft.Json documentation. \url{https://www.newtonsoft.com/json}
\end{enumerate}

\newpage
\appendix
\section{Приложение А. Текст программы}
В приложении приведены ключевые файлы проекта бухгалтерского приложения \textit{BuhUchet}.

\subsection{BuhUchet.csproj}
\lstinputlisting[language=XML]{BuhUchet.csproj}

\subsection{MainWindow.xaml}
\lstinputlisting[language=XML]{MainWindow.xaml}

\subsection{MainWindow.xaml.cs}
\lstinputlisting[language={[Sharp]C}]{MainWindow.xaml.cs}

\subsection{Models.cs}
\lstinputlisting[language={[Sharp]C}]{Models.cs}

\subsection{DataService.cs}
\lstinputlisting[language={[Sharp]C}]{DataService.cs}

\subsection{AutoCompleteService.cs}
\lstinputlisting[language={[Sharp]C}]{AutoCompleteService.cs}

\subsection{AppDataStore.cs}
\lstinputlisting[language={[Sharp]C}]{AppDataStore.cs}

\end{document}
